\begin{textsample}{NEG Dim 1 – unprofiled\_gpt – Score -10.00 – unprofiled\_gpt/t000414\_gpt.txt}  \label{ex:f1_neg_023}
I \textbf{just} think universities have got it \textbf{totally} the wrong way round.
They tend to shove the weakest, most limited lecturers in front of first‑year students, when actually that ’s when you need the very best people.
Those first encounters should be with really excellent, broad‑minded lecturers who can communicate to \textbf{anyone}, whether it ’s ten people or two hundred.
For me, a good lecturer should be able to \textbf{stand} there and \textbf{just} give, regardless of how many students are in the room.
It ’s a bit like a DJ : a \textbf{proper} DJ doesn't go, “ Oh, there are only fifteen people here, I won't \textbf{bother}. ” They still put \textbf{everything} into it, they still \textbf{try} to create something exciting and generous for whoever ’s turned up.
But a lot of lecturers \textbf{just} don't seem able to be generous with what they know.
It ’s as if something in them is blocked.
Their knowledge stays locked inside their \textbf{heads} rather than flowing outwards.
And that ’s such a waste, because students arrive at university \textbf{expecting} to be excited, and they ’re met with people who can't or won't really share what they ’ve got.

% matched lemmas: anyone, bother, everything, expect, head, just, proper, stand, totally, try
\end{textsample}
